\chapter{Ruch Browna w wielu wymiarach i dalsze własności trajektorii}
\label{ch:brown-multi}

\section{Dlaczego potrzebujemy procesu wielowymiarowego?}
W ruchu jednowymiarowym możemy badać przekroczenie bariery przez
cząstkę. Ruch w przestrzeni fizycznej wymaga jednak opisu wielu
współrzędnych naraz. Pytania zmieniają się wtedy jakościowo:
jak długo cząstka pozostaje w kuli? Czy dociera do wyróżnionego
punktu? Jak zmieniają się wyniki po obróceniu układu współrzędnych?
Zacznijmy od konstrukcji, w której odpowiedź na ostatnie pytanie
wynika ze szczególnej postaci macierzy kowariancji.

\begin{definicja}[Standardowy wielowymiarowy ruch Browna]
Niech $d\ge1$ będzie liczbą całkowitą. Proces $B=(B_t)_{t\ge0}$ o wartościach
w $\R^d$, z prawie na pewno ciągłymi trajektoriami, nazywamy
\emph{standardowym $d$-wymiarowym ruchem Browna}, jeżeli
$B_0=0$ prawie na pewno, ma niezależne przyrosty oraz
$B_t-B_s\sim\Normal_d(0,(t-s)I_d)$ dla wszystkich $0\le s<t$.
Jeżeli $x\in\R^d$, proces $X_t=x+B_t$ nazywamy ruchem Browna
rozpoczynającym się w $x$; jego prawo i wartość oczekiwaną oznaczamy
odpowiednio przez $\Prob_x$ oraz $\E_x$.
Gdy zadano filtrację $(\calF_t)_{t\ge0}$, wymagamy
dodatkowo adaptacji oraz niezależności przyrostu
$B_t-B_s$ od $\calF_s$ dla $s<t$.
W dalszym ciągu, o ile nie zaznaczono inaczej, wybieramy
wersję o wszystkich trajektoriach ciągłych i $B_0=0$
punktowo oraz jej filtrację naturalną.
\end{definicja}

Ruch Browna względem swojej filtracji naturalnej
spełnia warunek niezależności przyszłego przyrostu:
skończony wektor obserwacji do chwili $s$ jest
prawie na pewno funkcją przyrostów do tej chwili.
Niezależność przyrostu $B_t-B_s$ od zdarzeń
cylindrycznych rozszerzamy argumentem $\pi$--$\lambda$
na $\calF_s=\sigma(B_v:0\le v\le s)$.

\begin{twierdzenie}[Konstrukcja i izotropia]
Dla każdego $d\ge1$ istnieje standardowy $d$-wymiarowy ruch
Browna $B$. Jego współrzędne są niezależnymi standardowymi
procesami Wienera, a dla $1\le i,j\le d$ i $s,t\ge0$
\[
\E B_t^{(i)}=0,\qquad
\Cov(B_s^{(i)},B_t^{(j)})=\delta_{ij}\min(s,t).
\]
Jeżeli $U\in\R^{d\times d}$ jest macierzą ortogonalną,
to proces $(UB_t)_{t\ge0}$ ma to samo prawo na
$C_0([0,T];\R^d)$ co $B$, dla każdego $T>0$.
\end{twierdzenie}
\begin{proof}
Na przestrzeni iloczynowej $d$ kopii przestrzeni probabilistycznej
z ruchem Browna weźmy niezależne procesy $W^{(1)},\ldots,W^{(d)}$
i połóżmy $B_t=(W_t^{(1)},\ldots,W_t^{(d)})$.
Wektor każdego przyrostu ma niezależne współrzędne o rozkładzie
$\Normal(0,t-s)$, czyli ma rozkład $\Normal_d(0,(t-s)I_d)$.
Współrzędne na różnych rozłącznych przedziałach także są niezależne;
można to zobaczyć, grupując skończone wektory przyrostów po
współrzędnych. Ciągłość wynika z przecięcia $d$ zdarzeń pełnej miary.

Dla dowolnego procesu spełniającego definicję wektor
wartości w skończonym zbiorze chwil jest prawie na pewno
liniowym obrazem wektora niezależnych gaussowskich przyrostów
(zaczynamy od chwili zero). Jest zatem gaussowski.
Dla $s\le t$ niezależność $B_t-B_s$ od $B_s$ daje
$\Cov(B_s^{(i)},B_t^{(j)})=\Cov(B_s^{(i)},B_s^{(j)})
=s\delta_{ij}$; średnie są zerowe. Dla dowolnych skończonych
rodzin obserwacji poszczególnych współrzędnych macierz
kowariancji jest blokowo diagonalna względem współrzędnych,
więc odpowiednie bloki są łącznie niezależne. Iterując
argument $\pi$--$\lambda$ po $d$ współrzędnych, otrzymujemy
niezależność $\sigma$-ciał generowanych przez całe procesy
współrzędnych, nie tylko przez ich wartości w jednej chwili.
Każdy proces współrzędnej jest standardowym procesem Wienera.

Ponieważ $U((t-s)I_d)U^\top=(t-s)I_d$, przyrosty procesu $UB$
mają te same rozkłady co przyrosty $B$ i pozostają niezależne.
Stąd równość rozkładów skończeniewymiarowych, a przez generowanie borelowskiego $\sigma$-ciała
przestrzeni $C_0([0,T];\R^d)$ przez ewaluacje w przeliczalnym
gęstym zbiorze chwil także równość praw trajektorii.
Wynika to z jednowymiarowego argumentu z rozdziału~\ref{ch:wiener}
i identyfikacji przestrzeni trajektorii wektorowych z iloczynem
$d$ przestrzeni trajektorii rzeczywistych.
\end{proof}

\begin{stwierdzenie}[Równoważna charakteryzacja gaussowska]
Rzeczywisty proces wektorowy $B$ o prawie na pewno ciągłych
trajektoriach jest standardowym ruchem Browna w $\R^d$
wtedy i tylko wtedy, gdy jest wycentrowanym procesem
gaussowskim o macierzach kowariancji
$\Cov(B_s,B_t)=\min(s,t)I_d$.
\end{stwierdzenie}
\begin{proof}
Konieczność wykazano w poprzednim twierdzeniu.
Odwrotnie, z zerowej średniej i kowariancji w chwili zero
wynika $B_0=0$ p.n. Przyrost $B_t-B_s$ jest gaussowski
o średniej zero i kowariancji $(t-s)I_d$.
Dla $0\le a<b\le c<e$ kowariancja przyrostów
$B_b-B_a$ i $B_e-B_c$ jest równa
$(b-b-a+a)I_d=0$. Dowolny skończony wektor takich
przyrostów jest gaussowski o blokowo diagonalnej
macierzy kowariancji, co daje ich łączną niezależność.
\end{proof}

\begin{uwaga}[Niezależność współrzędnych a niezależność położenia]
Współrzędne $B_t^{(1)},\ldots,B_t^{(d)}$ są niezależne w ustalonej
chwili $t$. Nie wynika stąd, że położenia $B_s$ i $B_t$ są niezależne
dla $0<s<t$: przeciwnie, mają wspólną przeszłość, a ich kowariancja
wynosi $sI_d$. Przy $s=0$ położenie początkowe jest stałe p.n. Niezależne są \emph{przyrosty} na rozłącznych
przedziałach czasu.
\end{uwaga}

\begin{twierdzenie}[Mocna własność Markowa w wielu wymiarach]
\label{thm:brown-multi-strong}
Niech $B$ będzie standardowym ruchem Browna względem
filtracji $(\calF_t)_{t\ge0}$, o wszystkich trajektoriach
ciągłych, a $\tau$ czasem zatrzymania skończonym prawie
na pewno. Wówczas proces
$(B_{\tau+v}-B_\tau)_{v\ge0}$ (z trajektorią zerową
przy $\tau=\infty$) jest standardowym $d$-wymiarowym
ruchem Browna niezależnym od $\calF_\tau$.
\end{twierdzenie}
\begin{proof}
Dowód jest wektorową wersją dowodu
twierdzenia~\ref{thm:wiener-strong-markov}.
Dla $h_n=2^{-n}$ zaokrąglamy $\tau$ w górę do
$\tau_n=h_n\lceil\tau/h_n\rceil$ (z wartością
$\infty$ przy $\tau=\infty$). Są to czasy zatrzymania,
$\tau_n\downarrow\tau$, a $\calF_\tau\subseteq\calF_{\tau_n}$.
Obserwacje w losowych chwilach są mierzalne na mocy
twierdzeń~\ref{thm:prog} i~\ref{thm:sampled}, ponieważ
$\tau+v$ i $\tau_n+v$ są czasami zatrzymania dla $v\ge0$.
Przy nieskończonej chwili przyjmujemy wartość zero
dla wszystkich takich obserwacji.
Dla $A\in\calF_\tau$, chwil $0\le v_1<\cdots<v_m$
i wektorów $u_1,\ldots,u_m\in\R^d$ warunkujemy na
zdarzeniach $A\cap\{\tau_n=kh_n\}\in\calF_{kh_n}$.
Kolejne przyrosty po $kh_n$ są łącznie niezależne
od $\calF_{kh_n}$ przez iterowane warunkowanie
względem chwil ich rozpoczęcia. Z gaussowskości
i wzoru na kowariancję dostajemy po zsumowaniu
po $k\ge0$
\[
\E\left[\chi_A\exp\left(
i\sum_{j=1}^m\langle u_j,B_{\tau_n+v_j}-B_{\tau_n}\rangle
\right)\right]
=\Prob(A)\exp\left(
-\frac12\sum_{j,l=1}^m\min(v_j,v_l)\langle u_j,u_l\rangle
\right).
\]
Sumowanie jest dozwolone przez ograniczenie modułu
funkcji wykładniczej przez $1$ i fakt, że zdarzenie
$\{\tau_n=\infty\}$ ma prawdopodobieństwo zero.
Z ciągłości trajektorii i zbieżności zmajoryzowanej
ta sama równość zachodzi z $\tau$ w miejsce $\tau_n$.
Dla $A=\Omega$ identyfikuje ona rozkłady
skończeniewymiarowe ruchu Browna, a dla dowolnego
$A$ dowodzi ich niezależności od $\calF_\tau$
przez jednoznaczność funkcji charakterystycznych.
Ciągłość przyszłych trajektorii i mierzalność
poszczególnych obserwacji zapewniają mierzalność
odwzorowania trajektorii, ponieważ jego borelowskie
$\sigma$-ciało jest generowane przez ewaluacje
w chwilach wymiernych. Argument $\pi$--$\lambda$
dla cylindrów rozszerza niezależność na całe to
$\sigma$-ciało i kończy dowód.
\end{proof}

\begin{twierdzenie}[Generator i wielowymiarowe równanie ciepła]
\label{thm:brown-multi-generator}
Niech $B$ będzie standardowym ruchem Browna w $\R^d$.
Dla $t\ge0$ oraz ograniczonej funkcji borelowskiej
$f:\R^d\to\R$ połóżmy $P_tf(x)=\E f(x+B_t)$.
Operatory te tworzą półgrupę przejścia: $P_0f=f$,
$P_sP_t=P_{s+t}$, a względem filtracji naturalnej
dla $X=x+B$ zachodzi
$\E[f(X_{s+t})\mid\calF_s]=(P_tf)(X_s)$ p.n.
Jeżeli $f\in C_b^3(\R^d)$ (wszystkie pochodne do rzędu trzeciego
są ciągłe i ograniczone, włącznie z samą funkcją), to
\[
\lim_{t\downarrow0}\left\|\frac{P_tf-f}{t}-\frac12\Delta f
\right\|_\infty=0.
\]
Jeżeli $f\in C_b(\R^d)$, to $u(t,x)=P_tf(x)$ jest klasy
$C^{1,2}$ na $(0,\infty)\times\R^d$, spełnia
$\partial_tu=\frac12\Delta u$ oraz $u(t,x)\to f(x)$ dla $t\downarrow0$.
\end{twierdzenie}
\begin{proof}
Mierzalność $P_tf$ oraz wzór na warunkowanie
wynikają z twierdzenia Fubiniego i niezależności
$B_{s+t}-B_s$ od $\calF_s$, tak jak w rozdziale
\ref{ch:wiener}. Własność półgrupy wynika z faktu,
że suma niezależnych wektorów o rozkładach
$\Normal_d(0,sI_d)$ i $\Normal_d(0,tI_d)$
ma rozkład $\Normal_d(0,(s+t)I_d)$.

Z wielowymiarowego wzoru Taylora z resztą, ograniczeniem
pochodnych trzeciego rzędu i normą euklidesową dostajemy
$f(x+h)=f(x)+\langle\nabla f(x),h\rangle+
\frac12 h^\top D^2f(x)h+R_x(h)$, gdzie
$|R_x(h)|\le C_d\max_{|\alpha|=3}\|\partial^\alpha f\|_\infty|h|^3$.
Ponieważ $\E B_t=0$,
$\E[B_t^{(i)}B_t^{(j)}]=t\delta_{ij}$ i
$\E|B_t|^3=t^{3/2}\E|B_1|^3$, otrzymujemy
$P_tf(x)-f(x)=\frac{t}{2}\sum_i\partial_{ii}f(x)+O(t^{3/2})$.
Stała w reszcie jest niezależna od $x$, a trzeci moment
normy $B_1$ jest skończony, ponieważ
$|B_1|^3\le d^{1/2}\sum_i|B_1^{(i)}|^3$.
To dowodzi granicy w normie supremum po podzieleniu przez $t$.
Dla $t>0$ rozkład $B_t$ ma gęstość
$p_t^{(d)}(z)=(2\pi t)^{-d/2}\exp(-|z|^2/(2t))$.
Na każdym pasie czasowym $0<\delta\le t\le R$ mamy
\[
\partial_t p_t^{(d)}(z)=
\left(-\frac d{2t}+\frac{|z|^2}{2t^2}\right)p_t^{(d)}(z),
\qquad
\partial_{ij}p_t^{(d)}(z)=
\left(\frac{z_i z_j}{t^2}-\frac{\delta_{ij}}{t}\right)p_t^{(d)}(z).
\]
Ponadto $\partial_i p_t^{(d)}(z)=-(z_i/t)p_t^{(d)}(z)$.
Piszemy $u(t,x)=\int_{\R^d} f(y)p_t^{(d)}(y-x)\,
\mathrm d\lambda_d(y)$, aby różniczkować jądro, a nie $f$.
Dla $|x-x_0|\le1$ mamy $|y-x|^2\ge |y-x_0|^2/2-1$.
Zatem wszystkie wskazane pochodne pod znakiem całki
mają wspólną całkowalną majorantę, będącą stałą razy
$(1+|y-x_0|^2)\exp(-|y-x_0|^2/(4R))$.
Różniczkowanie pod znakiem całki i twierdzenie o zbieżności
zmajoryzowanej dają ciągłość pochodnych oraz
$\partial_tu=\frac12\Delta u$, gdyż
$\partial_t p_t^{(d)}=\frac12\Delta p_t^{(d)}$.
Zbieżność do warunku początkowego wynika z ciągłości $f$ w $x$,
ograniczenia $f$ i faktu $\E|B_t|^2=dt\to0$, dokładnie jak
w dowodzie jednowymiarowym.
\end{proof}

\begin{przyklad}[Ruch Browna z anizotropową kowariancją]
Niech $A\in\R^{m\times d}$, $b\in\R^m$ i $X_t=b+AB_t$.
Wówczas $X_t-X_s\sim\Normal_m(0,(t-s)AA^\top)$ i
$\E X_t=b$. Jeżeli $AA^\top$ jest dodatnio określona, ale nie jest wielokrotnością
macierzy jednostkowej, kierunki nie są równoważne: powierzchnie
jednakowej gęstości są elipsoidami. Dla
$f\in C_b^3(\R^m)$ generator ma postać
$Lf=\frac12\sum_{i,j}(AA^\top)_{ij}\partial_{ij}f$.
Wzór wynika z tego samego rachunku Taylora, po zastąpieniu
macierzy kowariancji $tI_m$ przez $tAA^\top$, dla półgrupy
$P_t^Af(x)=\E f(x+AB_t)$. Przy osobliwej macierzy
$AA^\top$ mamy $\operatorname{rank} A<m$, ponieważ
$\ker(AA^\top)=\ker A^\top$. Proces startujący w $b$
porusza się w $b+\operatorname{im}A$, podprzestrzeni
afinicznej o wymiarze mniejszym od $m$. Jej miara
Lebesgue'a jest zerowa, więc rozkład $X_t$ jest singularny
względem $\lambda_m$.
\end{przyklad}


\begin{zadanie}[Anizotropowa dyfuzja w płaszczyźnie]
Niech $B=(B^{(1)},B^{(2)})$ będzie standardowym ruchem Browna
w $\R^2$ i niech $X_t=(2B_t^{(1)},B_t^{(2)})$.
Wyznaczyć rozkład $X_t$, generator procesu i prawdopodobieństwo
$\Prob(|X_t^{(1)}|>2a)$ dla $a,t>0$.

\emph{Rozwiązanie.} Macierz $A=\operatorname{diag}(2,1)$ daje
$AA^\top=\operatorname{diag}(4,1)$. Zatem
$X_t\sim\Normal_2(0,t\operatorname{diag}(4,1))$,
a jego gęstość wynosi
$ (4\pi t)^{-1}\exp(-x_1^2/(8t)-x_2^2/(2t))$.
Generator ma postać $Lf=2\partial_{11}f+\tfrac12\partial_{22}f$.
Wreszcie
$\Prob(|X_t^{(1)}|>2a)=\Prob(|B_t^{(1)}|>a)
=2(1-\Phi(a/\sqrt t))$.
Współczynnik dyfuzji w pierwszym kierunku jest cztery razy
większy niż w drugim; nie można go odczytać z samej
wartości oczekiwanej $\E X_t=0$.
\end{zadanie}

\section{Czas wyjścia z kuli i martyngał kwadratu promienia}
\begin{twierdzenie}[Średni czas wyjścia z kuli]
Niech $d\ge1$ będzie liczbą całkowitą, $R>0$,
$x\in\R^d$ i $|x|<R$.
Dla standardowego $d$-wymiarowego ruchu Browna $B$ połóżmy
$X_t=x+B_t$ oraz $\tau=\inf\{t\ge0:|X_t|\ge R\}$.
Wtedy $\tau<\infty$ prawie na pewno i
\[
\E\tau=\frac{R^2-|x|^2}{d}.
\]
\end{twierdzenie}
\begin{proof}
Pracujemy z ciągłą wersją $B$ zaczynającą się w zerze
i jej filtracją naturalną. Niezależność przyrostu od
tej filtracji wynika ze zdarzeń cylindrycznych i argumentu
$\pi$--$\lambda$, tak jak w przypadku jednowymiarowym.
Czas $\tau$ jest czasem zatrzymania na mocy
twierdzenia~\ref{thm:hitting-closed}, zastosowanego do
domkniętego zbioru $\{y:|y|\ge R\}$.
Dla $s<t$ i $Z=B_t-B_s$, niezależnego od przeszłości,
$\E Z=0$ i $\E|Z|^2=d(t-s)$. Rozwinięcie kwadratu daje
$\E[|X_t|^2\mid\calF_s]=|X_s|^2+d(t-s)$.
Zatem $M_t=|X_t|^2-dt$ jest całkowalnym martyngałem
o ciągłych trajektoriach. Dla każdego całkowitego $n\ge1$
twierdzenie o próbkowaniu opcjonalnym
(twierdzenie~\ref{thm:optional-cont}) daje
\[
\E|X_{\tau\wedge n}|^2-|x|^2=d\E(\tau\wedge n).
\]
Z ciągłości trajektorii zawsze $|X_{\tau\wedge n}|\le R$:
przed $\tau$ proces jest w kuli, a w chwili $\tau<\infty$
osiąga jej sferę. Stąd $d\E(\tau\wedge n)\le R^2-|x|^2$.
Zbieżność monotoniczna daje $\E\tau\le (R^2-|x|^2)/d<\infty$,
co wymusza $\tau<\infty$ p.n. Wtedy
$|X_{\tau\wedge n}|^2\to|X_\tau|^2=R^2$ p.n., a
zmienne są ograniczone przez $R^2$. Zbieżność zmajoryzowana
i monotoniczna po obu stronach poprzedniej tożsamości dają
$R^2-|x|^2=d\E\tau$.
\end{proof}

\begin{zadanie}[Wyjście z dysku]
Rozważmy standardowy dwuwymiarowy ruch Browna, którego
położenie początkowe ma odległość $r<R$ od środka dysku
promienia $R$. Ile średnio trzeba czekać na pierwsze osiągnięcie
okręgu?

\emph{Rozwiązanie.} Z twierdzenia dla $d=2$ dostajemy
$\E\tau=(R^2-r^2)/2$. Gdy $r=0$, średni czas to $R^2/2$.
Zależność od $r$ jest kwadratowa, mimo że typowe wychylenie
po czasie $t$ ma rząd $\sqrt t$.
\end{zadanie}

\section{Trafianie w kule i punkty: zależność od wymiaru}

Aby zbadać trafianie w małe zbiory, potrzebujemy
martyngałów związanych z funkcjami harmonicznymi.
Wyprowadzimy je bez stosowania wzoru Itô, który
zostanie wprowadzony w tomie III.

\begin{lemat}[Martyngał związany z generatorem]
\label{lem:brown-generator-martingale}
Dla $f\in C_c^2(\R^d)$ i $X_t=x+B_t$ proces
\[
M_t^f=f(X_t)-f(x)-\frac12\int_0^t
\Delta f(X_v)\,\mathrm d\lambda_1(v)
\]
jest ciągłym martyngałem względem filtracji naturalnej.
\end{lemat}
\begin{proof}
Najpierw przyjmijmy $f\in C_c^\infty(\R^d)$.
Oznaczmy $g=\frac12\Delta f$. Dla $0\le s<t$
weźmy podział $v_j=s+jh$, $h=(t-s)/n$.
Niezależność przyrostów od przeszłości daje
\[
\E[f(X_{v_{j+1}})-f(X_{v_j})\mid\calF_{v_j}]
=(P_h f-f)(X_{v_j}).
\]
Z twierdzenia~\ref{thm:brown-multi-generator}
mamy $P_h f-f=hg+h r_h$, gdzie
$\norm{r_h}_\infty\to0$. Stosując własność wieży
i sumując, otrzymujemy
\[
\E[f(X_t)-f(X_s)\mid\calF_s]
=\E\left[\sum_{j=0}^{n-1}h g(X_{v_j})
+\sum_{j=0}^{n-1}h r_h(X_{v_j})\mid\calF_s\right].
\]
Druga suma jest ograniczona w module przez
$(t-s)\norm{r_h}_\infty\to0$. Pierwsza zbiega
punktowo do całki po $[s,t]$, a ponieważ $g$
jest ograniczona, również w $L^1$.
Kontraktywność warunkowej wartości oczekiwanej
w $L^1$ pozwala przejść do granicy. Otrzymujemy
$\E[M_t^f-M_s^f\mid\calF_s]=0$.
Adaptacja całki wynika z aproksymacji sumami
Riemanna, a jej ciągłość i całkowalność
z ciągłości trajektorii i ograniczoności $g$.

Dla ogólnego $f\in C_c^2(\R^d)$ weźmy sploty
$f_\varepsilon=f*\varphi_\varepsilon$ z nieujemną
gładką funkcją aproksymującą jedynkę, o całce $1$
i nośniku w kuli promienia $\varepsilon$.
Funkcje $f_\varepsilon$ są gładkie o zwartym nośniku,
a $f_\varepsilon\to f$ oraz
$\Delta f_\varepsilon\to\Delta f$ jednostajnie:
wynika to z jednostajnej ciągłości $f$ i jej drugich
pochodnych oraz różniczkowania splotu.
Dla ustalonego $T$ różnica procesów $M^{f_\varepsilon}$
i $M^f$ jest na $[0,T]$ ograniczona w module przez
$2\norm{f_\varepsilon-f}_\infty+
(T/2)\norm{\Delta f_\varepsilon-\Delta f}_\infty$.
Przechodząc do granicy w tożsamościach
martyngałowych w $L^1$, dostajemy tezę.
\end{proof}

\begin{twierdzenie}[Rozkład miejsca wyjścia z kuli ze środka]
\label{thm:brown-exit-sphere}
Niech $B$ będzie standardowym ruchem Browna w $\R^d$,
$R>0$ i $\sigma_R=\inf\{t\ge0:|B_t|\ge R\}$.
Zmienna $B_{\sigma_R}$ ma jednostajny rozkład na
sferze $S_R=\{y\in\R^d:|y|=R\}$, to znaczy
\[
\Prob(B_{\sigma_R}\in D)=
\frac{\mathcal H^{d-1}(D\cap S_R)}
{\mathcal H^{d-1}(S_R)}
\]
dla każdego zbioru borelowskiego $D\subseteq\R^d$.
W wymiarze $d=1$ miara $\mathcal H^0$ jest miarą
liczącą, a obydwa końce przedziału mają masę $1/2$.
Ponadto czas wyjścia $\sigma_R$ i miejsce wyjścia
$B_{\sigma_R}$ są niezależne.
\end{twierdzenie}
\begin{proof}
Czas $\sigma_R$ jest skończony prawie na pewno
na mocy twierdzenia o średnim czasie wyjścia.
Mierzalność $B_{\sigma_R}$ wynika z twierdzenia
\ref{thm:sampled}; na zbiorze zerowym nieskończonego
czasu można przyjąć dowolną stałą wartość.
Ciągłość trajektorii daje $|B_{\sigma_R}|=R$ p.n.
Dla dowolnej macierzy ortogonalnej $U$ czas wyjścia
procesu $UB$ jest taki sam jak czas wyjścia $B$.
Izotropia całego procesu daje więc
$\Law(UB_{\sigma_R})=\Law(B_{\sigma_R})$.
Wyjście jako funkcjonał trajektorii jest mierzalne:
czas wyjścia wyrażamy przez maksimum normy na
zwartych przedziałach, a miejsce wyjścia przez
granice ewaluacji w chwilach zaokrąglonych w górę.

Wykażemy jednoznaczność miary probabilistycznej
na $S_R$ niezmienniczej względem wszystkich
przekształceń ortogonalnych. Niech $Y$ ma taką miarę,
a $V=G/|G|$, gdzie $G\sim\Normal_d(0,I_d)$ jest
niezależny od $Y$ (w razie potrzeby na przestrzeni
iloczynowej) (na zdarzeniu $G=0$ przyjmujemy
dowolny wektor jednostkowy). Prawo $V$ jest
niezmiennicze względem przekształceń ortogonalnych.
Dla $u\in\R^d$ funkcja charakterystyczna $Y$
zależy tylko od $|u|$, więc
\[
\E e^{i\langle u,Y\rangle}
=\E e^{i|u|\langle V,Y\rangle}
=\E e^{i|u|R V_1}.
\]
W drugiej równości warunkujemy względem $Y$;
dla każdego $y\in S_R$ izotropia $V$ daje
$\langle V,y\rangle\sim R V_1$.
Prawa dowolnych dwóch miar niezmienniczych mają
zatem tę samą funkcję charakterystyczną.
Znormalizowana miara Hausdorffa na sferze
jest taką miarą, ponieważ izometrie zachowują
miarę Hausdorffa, a jej całkowita masa na sferze
jest skończona i dodatnia (dla $d\ge2$ wynika to
z lokalnych gładkich map sfery i wzoru na pole,
dla $d=1$ z liczenia punktów).
Jednoznaczność funkcji charakterystycznych dowodzi
pierwszej tezy.

Aby wykazać niezależność, ustalmy borelowski zbiór
$A\subseteq[0,\infty)$ o $\Prob(\sigma_R\in A)>0$.
Miara na sferze
\[
\nu_A(D)=
\frac{\Prob(\sigma_R\in A,\ B_{\sigma_R}\in D)}
{\Prob(\sigma_R\in A)}
\]
jest również niezmiennicza względem macierzy
ortogonalnych, gdyż obrót lub odbicie zachowuje
czas wyjścia, a przekształca położenie zgodnie z $U$.
Udowodniona jednoznaczność miary niezmienniczej
daje $\nu_A=\Law(B_{\sigma_R})$.
Dla $A$ o prawdopodobieństwie zero faktoryzacja
jest automatyczna. Otrzymujemy niezależność.
\end{proof}

\begin{wniosek}[Własność wartości średniej funkcji harmonicznej]
Niech $u$ będzie funkcją klasy $C^2$ na otwartym
otoczeniu domkniętej kuli $\{y:|y-x|\le R\}$,
harmoniczną w jej wnętrzu, czyli $\Delta u=0$.
Wtedy
\[
u(x)=
\frac{1}{\mathcal H^{d-1}(S_R)}
\int_{S_R}u(x+y)\,\mathrm d\mathcal H^{d-1}(y).
\]
\end{wniosek}
\begin{proof}
Wybieramy gładką funkcję odcinającą równą $1$
w otoczeniu domkniętej kuli, o zwartym nośniku
w obszarze definicji $u$. Po pomnożeniu przez $u$
i przedłużeniu przez zero otrzymujemy funkcję
$f\in C_c^2(\R^d)$ równą $u$ na kuli.
Niech $X_t=x+B_t$ i $\tau=\sigma_R$.
Z lematu~\ref{lem:brown-generator-martingale}
oraz twierdzenia~\ref{thm:optional-cont} wynika
$\E f(X_{\tau\wedge n})=f(x)$, ponieważ
$\Delta f(X_s)=0$ dla $s<\tau$.
Czas $\tau$ jest skończony p.n., a $f$ jest ograniczona.
Przechodząc do granicy przez twierdzenie
o zbieżności zmajoryzowanej, dostajemy
$u(x)=\E u(X_\tau)$. Twierdzenie
\ref{thm:brown-exit-sphere} daje wskazaną całkę.
\end{proof}

\begin{twierdzenie}[Trafienie w sferę wewnętrzną przed zewnętrzną]
\label{thm:brown-annulus}
Niech $0<r<|x|<R$ i $X_t=x+B_t$. Oznaczmy
$\tau_r=\inf\{t\ge0:|X_t|\le r\}$ oraz
$\sigma_R=\inf\{t\ge0:|X_t|\ge R\}$.
Wtedy
\[
\Prob_x(\tau_r<\sigma_R)=
\begin{cases}
\dfrac{R-|x|}{R-r},&d=1,\\[6pt]
\dfrac{\log(R/|x|)}{\log(R/r)},&d=2,\\[6pt]
\dfrac{|x|^{2-d}-R^{2-d}}{r^{2-d}-R^{2-d}},&d\ge3.
\end{cases}
\]
\end{twierdzenie}
\begin{proof}
Dla radialnej funkcji $u(y)=h(|y|)$ poza zerem
bezpośrednie różniczkowanie daje
$\Delta u(y)=h''(|y|)+(d-1)h'(|y|)/|y|$.
W istocie $\partial_i u=h'(\rho)y_i/\rho$,
gdzie $\rho=|y|$, oraz
\[
\partial_{ii}u(y)=h''(\rho)\frac{y_i^2}{\rho^2}
+h'(\rho)\left(\frac1\rho-\frac{y_i^2}{\rho^3}\right).
\]
Zsumowanie po $i$ daje wskazany wzór. Wybieramy $h(\rho)=\rho$
dla $d=1$, $h(\rho)=\log\rho$ dla $d=2$ i
$h(\rho)=\rho^{2-d}$ dla $d\ge3$.
W każdym przypadku $\Delta u=0$ poza zerem.
Mnożąc $u$ przez gładką funkcję odcinającą,
równą $1$ w otoczeniu domkniętego pierścienia
$\{r\le |y|\le R\}$, o nośniku zwartym poza zerem,
otrzymujemy $f\in C_c^\infty(\R^d)$ zgodną z $u$
na tym otoczeniu.

Czas $\eta=\tau_r\wedge\sigma_R$ jest czasem
zatrzymania na mocy twierdzenia~\ref{thm:hitting-closed}
i jest skończony p.n., gdyż $\eta\le\sigma_R$,
a skończoność $\sigma_R$ wynika z twierdzenia
o średnim czasie wyjścia z kuli.
Przed $\eta$ mamy $\Delta f(X_v)=0$.
Stosując lemat~\ref{lem:brown-generator-martingale}
oraz twierdzenie~\ref{thm:optional-cont} do
$\eta\wedge n$, otrzymujemy $\E_x f(X_{\eta\wedge n})=f(x)$.
Przez zbieżność zmajoryzowaną
$\E_x f(X_\eta)=f(x)$. Ciągłość trajektorii daje
$|X_\eta|\in\{r,R\}$; oba czasy nie mogą być
jednocześnie równe $\eta<\infty$.
Dla $p=\Prob_x(\tau_r<\sigma_R)$ otrzymujemy
$h(|x|)=p h(r)+(1-p)h(R)$.
Rozwiązując równanie względem $p$, dostajemy wszystkie
trzy wzory.
\end{proof}

\begin{wniosek}[Prawdopodobieństwo trafienia w domkniętą kulę]
Dla $|x|>r>0$ zachodzi
\[
\Prob_x(\tau_r<\infty)=
\begin{cases}
1,&d=1,2,\\
(r/|x|)^{d-2},&d\ge3.
\end{cases}
\]
\end{wniosek}
\begin{proof}
Zdarzenia $\{\tau_r<\sigma_R\}$ rosną przy
$R\to\infty$ wzdłuż rosnącego ciągu promieni.
Ich suma jest równa $\{\tau_r<\infty\}$:
jeśli $\tau_r<\infty$, trajektoria na $[0,\tau_r]$
jest ograniczona i mieści się w dostatecznie dużej
kuli. Wystarczy zatem przejść do granicy we
wzorach twierdzenia~\ref{thm:brown-annulus}.
\end{proof}

\begin{twierdzenie}[Trafianie w pojedyncze punkty]
Niech $X_t=x+B_t$ i niech $y$ będzie ustalonym
punktem różnym od $x$.
Jeżeli $d=1$, to $X$ trafia w $y$ prawie na pewno.
Jeżeli $d\ge2$, to z prawdopodobieństwem $1$
nie trafia w ten punkt $y$ w żadnej chwili.
Ponadto przy $d\ge2$, startując w $y$, proces
prawie na pewno nie wraca do $y$ w żadnej chwili $t>0$.
\end{twierdzenie}
\begin{proof}
Przez przesunięcie możemy przyjąć $y=0$.
Dla $d=1$ wniosek wynika z rozkładu czasu
trafienia i symetrii w rozdziale~\ref{ch:wiener}.
Dla $d\ge2$ i ustalonego $R>|x|>0$
trafienie w zero przed wyjściem z kuli promienia $R$
wymusza trafienie w każdą kulę o promieniu
$0<r<|x|$ przed tym wyjściem. Z twierdzenia
\ref{thm:brown-annulus} prawdopodobieństwo
takiego zdarzenia jest nie większe niż
$\Prob_x(\tau_r<\sigma_R)\to0$ przy $r\downarrow0$.
Każde trafienie w zero w skończonej chwili
musiałoby zajść przed wyjściem z dostatecznie
dużej kuli, więc suma po całkowitych $R>|x|$
ma prawdopodobieństwo zero.

W przypadku powrotu do punktu startowego $x=y$,
dla każdej ustalonej chwili $q>0$
mamy $\Prob_y(X_q=y)=0$, ponieważ rozkład $X_q$
ma gęstość. Warunkując względem $\calF_q$ i
korzystając z niezależności przyszłego ruchu Browna
od przeszłości (twierdzenie~\ref{thm:brown-multi-strong}
dla $\tau=q$), wnioskujemy, że prawdopodobieństwo
trafienia w $y$ w chwili $t\ge q$ jest zero.
Zastosowane warunkowanie jest uzasadnione twierdzeniem
Fubiniego: dla $z\in\R^d$ i ciągłej trajektorii $w$
funkcja $\inf_{0\le v\le k}|z+w(v)-y|$ jest ciągła
w parze $(z,w)$ z normą supremum na $[0,k]$.
Zdarzenie trafienia w chwili $t\ge q$ jest zatem
borelowskie jako przeliczalna suma zdarzeń
zerowania takiego minimum dla $k=1,2,\ldots$.
Dla $z\ne y$ jego prawdopodobieństwo wynosi zero
na mocy pierwszej części twierdzenia, a stan $z=y$
ma w chwili $q$ prawdopodobieństwo zero. Suma po
$q=1/n$ wyklucza wszystkie powroty w chwilach $t>0$.
\end{proof}

\begin{twierdzenie}[Powracalność i przejściowość ruchu Browna]
Dla dowolnego ustalonego punktu startowego $x$
i $d=1,2$ z prawdopodobieństwem $1$ ruch Browna
odwiedza każdy niepusty zbiór otwarty w dowolnie
późnych chwilach. Dla $d\ge3$ zachodzi prawie na pewno
\[
\lim_{t\to\infty}|x+B_t|=\infty.
\]
\end{twierdzenie}
\begin{proof}
Dla $d=1,2$, dowolnej otwartej kuli $U$
i całkowitej chwili $n\ge1$ prawdopodobieństwo
wejścia do $U$ w chwili $t\ge n$ wynosi $1$:
warunkujemy względem $\calF_n$ i stosujemy wniosek
o trafieniu w domkniętą kulę zawartą w $U$
(z dowolnego punktu startowego). Jeśli stan należy
już do tej mniejszej kuli, trafienie jest natychmiastowe.
Mierzalność zdarzeń wejścia do otwartej kuli wynika
z ciągłości i wystarczy sprawdzać chwile wymierne.
Przecinamy zbiory pełnej miary po $n$ oraz
po przeliczalnej bazie kul o wymiernych środkach
i dodatnich wymiernych promieniach.

Dla $d\ge3$ ustalmy $r>0$ i wybierzmy rosnące
promienie $R_n>\max(r,|x|)$, $R_n\to\infty$.
Czasy $\sigma_{R_n}$ są skończone p.n. i
$|X_{\sigma_{R_n}}|=R_n$.
Mocna własność Markowa
(twierdzenie~\ref{thm:brown-multi-strong}) oraz
wzór na trafienie w kulę dają (na mocy niezależności
przyszłej trajektorii od $\calF_{\sigma_{R_n}}$,
warunkujemy po położeniu $X_{\sigma_{R_n}}$;
zdarzenie trafienia jest borelowskie, gdyż można
je opisać minimum odległości na przedziałach
$[0,k]$, $k\ge1$)
\[
\Prob_x(\exists t\ge\sigma_{R_n}:|X_t|\le r)
=(r/R_n)^{d-2}\longrightarrow0.
\]
Zdarzenia po lewej są malejące (czasy wyjścia
rosną), zatem ich przecięcie ma miarę zero.
Prawie na pewno dla pewnego $n$ nie ma już
wizyt w kuli promienia $r$ po $\sigma_{R_n}$.
Przecinając zbiory pełnej miary dla $r=1,2,\ldots$,
otrzymujemy ucieczkę do nieskończoności.
\end{proof}

\begin{uwaga}[Powracalność w płaszczyźnie a trafianie w punkty]
Dwuwymiarowy ruch Browna prawie na pewno wchodzi
dowolnie późno do każdego otoczenia każdego punktu,
lecz ustalony punkt różny od punktu startowego
prawie na pewno nie zostaje trafiony.
Powracalność dotyczy otoczeń, nie pojedynczych punktów.
\end{uwaga}

\section{Maksimum ruchu Browna i mostu Browna}
Pełną zasadę odbicia i łączny rozkład maksimum
$M_T=\sup_{0\le s\le T}W_s$ oraz położenia $W_T$
udowodniono w twierdzeniu~\ref{thm:wiener-reflection}
i wniosku~\ref{cor:wiener-max-joint}.
Przypomnijmy ich konsekwencję: dla $a>0$, $T>0$ i
borelowskiego $A\subset(-\infty,a)$
\[
\Prob(M_T<a,\ W_T\in A)
=\int_A(p_T(y)-p_T(2a-y))\,\mathrm d\lambda_1(y),
\]
gdzie $p_T$ jest gęstością $\Normal(0,T)$.
Ponieważ $M_T$ ma prawo $|W_T|$, z momentu bezwzględnego
standardowej zmiennej normalnej otrzymujemy
$\E M_T=\sqrt{2T/\pi}$.
Wykorzystamy teraz te wyniki do analizy mostu Browna.

\begin{wniosek}[Maksimum mostu Browna]
Niech $W$ będzie standardowym procesem Wienera,
$T>0$ i $\beta_t=W_t-(t/T)W_T$ dla $0\le t\le T$.
Dla każdego $a>0$ zachodzi
\[
\Prob\left(\sup_{0\le t\le T}\beta_t\le a\right)
=1-\exp\left(-\frac{2a^2}{T}\right).
\]
\end{wniosek}
\begin{proof}
Z rozdziału~\ref{ch:gauss} wiemy, że $\beta$ jest
niezależny od $W_T$. Dla każdego $y\in\R$ proces
$\beta_t+(t/T)y$ ma rozkład warunkowy całej trajektorii
Wienera przy $W_T=y$; jest to \emph{regularna wersja}
rozkładu warunkowego, a nie warunkowanie na zdarzeniu o
prawdopodobieństwie dodatnim.

Niech $V=\sup_{0\le t\le T}\beta_t$ i dla $b>0$, $y<b$ połóżmy
$H(y,b)=\Prob(\sup_{0\le t\le T}(\beta_t+ty/T)<b)$.
Z twierdzenia~\ref{thm:wiener-reflection} i powyższego wzoru,
przez podzielenie gęstości
wspólnej zdarzenia braku przekroczenia i położenia końcowego
przez $p_T(y)>0$, otrzymujemy dla prawie każdego $y<b$
\[
H(y,b)=1-\frac{p_T(2b-y)}{p_T(y)}.
\]
Wyrażenie po prawej stronie ma dla $y\to0$ granicę
$g(b):=1-\exp(-2b^2/T)$, o ile $b>0$.

Ustalmy $a>0$ i $0<\varepsilon<a$. Jeśli $|y|<\min(\varepsilon,a-\varepsilon)$,
to oba progi $a-\varepsilon,a+\varepsilon$ są większe
od $y$, a dla każdej trajektorii $\sup_t|(t/T)y|\le\varepsilon$.
Dlatego mamy zawieranie zdarzeń prowadzące do nierówności
\[
H(y,a-\varepsilon)\le\Prob(V<a)
\le H(y,a+\varepsilon).
\]
Możemy wybrać ciąg $y_k\to0$ spoza sumy dwóch zbiorów
wyjątkowych miary Lebesgue'a zero, odpowiadających progom
$a-\varepsilon$ i $a+\varepsilon$. Podstawienie wzoru
na $H$ i przejście do granicy daje
$g(a-\varepsilon)\le\Prob(V<a)\le g(a+\varepsilon)$.
Ponieważ $g$ jest ciągła, po $\varepsilon\downarrow0$
otrzymujemy $\Prob(V<a)=g(a)$. Wreszcie
$\{V\le a\}=\bigcap_{n\ge1}\{V<a+1/n\}$;
ciągłość miary z góry i ciągłość $g$ dają
$\Prob(V\le a)=g(a)$. Dowód przy okazji wykazał,
że $\Prob(V=a)=0$ dla każdego $a>0$.
\end{proof}


\begin{uwaga}[Dlaczego przejście $y\to0$ nie jest tautologią?]
Dla ciągłej zmiennej $W_T$ zdarzenie $\{W_T=0\}$ ma
prawdopodobieństwo zero. Reguła Bayesa dla zdarzeń o dodatniej
mierze nie daje zatem rozkładu mostu przy zerowym punkcie
końcowym. Poprzedni dowód najpierw konstruuje \emph{konkretną}
wersję warunkowego prawa przy każdym $y$ przez rozkład
gaussowski, a następnie wykorzystuje ciągłość przesunięcia
trajektorii o funkcję $ty/T$. Wartość w $y=0$ jest tym
samym uzasadniona przez tę wersję, nie przez arbitralne
podstawienie do tożsamości obowiązującej tylko prawie wszędzie.
\end{uwaga}

\begin{zadanie}[Maksimum ruchu Browna w jednostce czasu]
Niech $W$ będzie standardowym procesem Wienera.
Obliczyć $\Prob(\sup_{0\le s\le1}W_s\ge1)$ oraz
$\E\sup_{0\le s\le1}W_s$.

\emph{Rozwiązanie.} Z zasady odbicia pierwsze
prawdopodobieństwo wynosi $2(1-\Phi(1))\approx0{,}3173$.
Wartość oczekiwana maksimum to $\sqrt{2/\pi}$.
Obserwujemy, że $\Prob(W_1\ge1)=1-\Phi(1)$ jest dwa razy
mniejsze: część trajektorii zdążyła osiągnąć próg,
a później wróciła poniżej niego.
\end{zadanie}


\begin{zadanie}[Warunkowy rozkład położenia mostu w dwóch wymiarach]
Niech $B$ będzie standardowym ruchem Browna w $\R^2$.
Wyznaczyć regularny rozkład warunkowy $B_1$ przy $B_4=y$,
gdzie $y=(2,-2)$, a następnie obliczyć prawdopodobieństwo,
że pierwsza współrzędna $B_1$ jest dodatnia.

\emph{Rozwiązanie.} Nie wolno dzielić przez $\Prob(B_4=y)=0$.
Korzystamy z rozkładu warunkowego wektorów gaussowskich.
Mamy $\Cov(B_1,B_4)=I_2$ oraz $\Var(B_4)=4I_2$.
Średnia warunkowa jest więc równa $I_2(4I_2)^{-1}y=y/4$,
a kowariancja wynosi
$I_2-I_2(4I_2)^{-1}I_2=\tfrac34 I_2$.
Zatem
\[
\Law(B_1\mid B_4=y)=\Normal_2\!\left(\frac{y}{4},\frac34 I_2\right).
\]
Dla $y=(2,-2)$ pierwsza współrzędna ma rozkład
$\Normal(1/2,3/4)$, skąd
$\Prob(B_1^{(1)}>0\mid B_4=y)=\Phi(1/\sqrt3)$.
Wynik opisuje \emph{wersję} warunkowania określoną
przez ciągłą funkcję parametru $y$.
\end{zadanie}

\begin{zadanie}[Warunkowanie położenia procesu Wienera w środku przedziału]
Niech $W$ będzie standardowym procesem Wienera. Wyznaczyć
regularny rozkład warunkowy $W_1$ przy $W_2=y$ oraz
obliczyć $\Prob(W_1>0\mid W_2=2)$ w wybranej wersji
rozkładu warunkowego.

\emph{Rozwiązanie.} Wektor $(W_1,W_2)$ jest gaussowski,
$\Var(W_1)=1$, $\Var(W_2)=2$ i
$\Cov(W_1,W_2)=1$. Wzór na warunkowanie wektorów
gaussowskich daje
\[
\Law(W_1\mid W_2=y)=\Normal(y/2,1/2).
\]
Nie dzielimy przez $\Prob(W_2=y)$, bo jest ono zerowe:
wykorzystujemy regularną wersję warunkowego rozkładu.
Przy $y=2$ otrzymujemy
\[
\Prob(W_1>0\mid W_2=2)=\Phi(\sqrt 2).
\]
Średnia warunkowa $y/2$ jest liniową interpolacją
między punktem początkowym $0$ i zaobserwowanym punktem
$y$; nie znaczy to jednak, że cała trajektoria mostu
jest deterministyczną prostą. Warunkowa wariancja jest
dodatnia i wynosi $1/2$.
\end{zadanie}

\section{Oscylacje drogi na dowolnie małych przedziałach}
\begin{twierdzenie}[Natychmiastowe oscylacje przy zadanej chwili]
Niech $W$ będzie standardowym procesem Wienera i niech
$t_0\ge0$ będzie ustaloną deterministyczną chwilą.
Wtedy z prawdopodobieństwem $1$ dla każdego $\varepsilon>0$
istnieją $s_+,s_-\in(t_0,t_0+\varepsilon)$ takie, że
$W_{s_+}>W_{t_0}$ oraz $W_{s_-}<W_{t_0}$.
\end{twierdzenie}
\begin{proof}
Proces $B_s=W_{t_0+s}-W_{t_0}$ jest standardowym
ruchem Browna. Dla każdego $r>0$ i $a>0$ zasada odbicia
daje
$\Prob(\sup_{0\le s\le r}B_s\ge a)=2(1-\Phi(a/\sqrt r))$.
Gdy $a\downarrow0$, prawdopodobieństwo to dąży do $1$.
Z ciągłości miary dla rosnących zdarzeń dostajemy
$\Prob(\sup_{0\le s\le r}B_s>0)=1$.
Ta sama własność dla $-B$ daje
$\Prob(\inf_{0\le s\le r}B_s<0)=1$.
Przecinamy zdarzenia pełnej miary dla $r=1/n$,
$n\in\N$. Z tego wynika teza dla wszystkich
$\varepsilon>0$, przez wybór $1/n<\varepsilon$.
\end{proof}

\begin{wniosek}[Natychmiastowe oscylacje po czasie zatrzymania]
Niech $W$ będzie procesem Wienera względem filtracji
$(\calF_t)_{t\ge0}$ o ciągłych trajektoriach i niech
$\tau$ będzie czasem zatrzymania skończonym p.n.
Wtedy prawie na pewno dla każdego $\varepsilon>0$
istnieją $u_+,u_-\in(0,\varepsilon)$ takie, że
\[
W_{\tau+u_+}>W_\tau,\qquad
W_{\tau+u_-}<W_\tau.
\]
\end{wniosek}
\begin{proof}
Twierdzenie~\ref{thm:wiener-strong-markov} daje
nowy proces Wienera $V_u=W_{\tau+u}-W_\tau$
na zdarzeniu skończoności $\tau$, z trajektorią
zerową na jego dopełnieniu. Stosujemy twierdzenie
o natychmiastowych oscylacjach do $V$ przy $t_0=0$.
Wspólny zbiór pełnej miary dla wszystkich
$\varepsilon>0$ otrzymujemy z przeliczalnej
rodziny przedziałów o długości $1/n$.
\end{proof}

\begin{uwaga}[Nie wolno zamienić ``każda ustalona chwila'' na ``wszystkie chwile'']
Udowodnione zdarzenie pełnej miary zależy od $t_0$.
Na ogół nie wolno przecinać nieprzeliczalnie wielu
zbiorów pełnej miary. Własność stwierdzona w twierdzeniu
nie zachodzi jednocześnie dla wszystkich chwil trajektorii:
na przykład maksimum na $[0,T]$ jest osiągane
w pewnej chwili $t_*<T$ prawie na pewno: proces odwrócony
$(W_T-W_{T-s})_{0\le s\le T}$ ma prawo Wienera
i przyjmuje ujemne wartości prawie na pewno, więc $W_T$
nie jest maksimum. Dla takiej chwili $t_*$ nie może
wystąpić większa wartość przed końcem przedziału. To ważna lekcja dotycząca kwantyfikatorów
w twierdzeniach o trajektoriach.
W szczególności chwila pierwszego osiągnięcia
maksimum na całym $[0,T]$ nie jest czasem zatrzymania
względem naturalnej filtracji: gdyby była, poprzedni
wniosek wymuszałby przekroczenie tego maksimum
przed $T$ na zdarzeniu $\{t_*<T\}$.
Mierzalność pierwszej chwili maksimum jako funkcjonału
całej trajektorii wynika z
$\{t_*\le s\}=\{\max_{[0,s]}W=\max_{[0,T]}W\}$
dla $0\le s\le T$.
\end{uwaga}

\section{Wielowymiarowa wariacja kwadratowa}
\begin{twierdzenie}[Macierz wariacji kwadratowej ruchu Browna]
\label{thm:brown-multi-bracket}
Niech $T>0$, $B$ będzie standardowym ruchem Browna
w $\R^d$ i niech $\pi=\{0=t_0<\cdots<t_m=T\}$.
Dla $t\in[0,T]$ oznaczmy
$D_k(t)=B_{t\wedge t_{k+1}}-B_{t\wedge t_k}$ oraz
\[
Q^\pi_t=\sum_{k=0}^{m-1}D_k(t)D_k(t)^\top.
\]
Dla dowolnego deterministycznego ciągu podziałów
$\pi_n$ o średnicach dążących do zera zachodzi
\[
\sup_{0\le t\le T}\norm{Q^{\pi_n}_t-tI_d}
\longrightarrow0
\quad\text{według prawdopodobieństwa},
\]
gdzie można użyć dowolnej normy macierzowej.
Jeżeli ponadto $\sum_n|\pi_n|<\infty$, zbieżność
jest prawie pewna.
\end{twierdzenie}
\begin{proof}
Dla ustalonej współrzędnej stosujemy jednowymiarowy
rachunek wariacji kwadratowej z rozdziału~\ref{ch:wiener}.
Dokładniej, dla dowolnego podziału, $h_k=t_{k+1}-t_k$,
zmienne $Y_k=(B^{(i)}_{t_{k+1}}-B^{(i)}_{t_k})^2-h_k$
są niezależne, wycentrowane i mają wariancje $2h_k^2$.
Nierówność maksymalna Kołmogorowa daje
\[
\Prob\left(\max_{j\le m}\left|\sum_{k<j}Y_k\right|
>\varepsilon\right)
\le \frac{2T|\pi|}{\varepsilon^2}.
\]
Między węzłami błąd względem lewego węzła
nie przekracza kwadratu modułu ciągłości trajektorii
na skali $|\pi|$ powiększonego o $|\pi|$.
Zatem mamy jednostajną zbieżność według
prawdopodobieństwa, a przy sumowalnych średnicach
także prawie pewną (Borel--Cantelli dla
$\varepsilon=1/l$, $l\ge1$).
Dla $i\ne j$ proces
$W=(B^{(i)}+B^{(j)})/\sqrt2$ jest standardowym
procesem Wienera na mocy charakterystyki gaussowskiej.
Z tożsamości polaryzacyjnej
\[
Q^\pi_{t,ij}
=V_t^\pi(W)-\tfrac12V_t^\pi(B^{(i)})
-\tfrac12V_t^\pi(B^{(j)}),
\]
gdzie $V^\pi$ oznacza jednowymiarową sumę kwadratów,
wynika zbieżność jednostajna do $t-t/2-t/2=0$.
Na przekątnej granica wynosi $t$.
Ponieważ współrzędnych jest skończenie wiele,
otrzymujemy tezę w dowolnej normie macierzowej
przez równoważność norm. W rozdziale
\ref{ch:quadratic-cont} powrócimy do tego pojęcia
dla ogólnych ciągłych martyngałów.
\end{proof}

\section{Rozkład promienia i oszacowania ryzyka wyjścia}
\begin{twierdzenie}[Rozkład odległości od początku]
Niech $B$ będzie standardowym ruchem Browna w $\R^d$
i $t>0$. Wtedy $|B_t|^2/t$ ma rozkład $\chi^2_d$,
a w szczególności
$\E|B_t|^2=dt$ i $\Var(|B_t|^2)=2dt^2$.
\end{twierdzenie}
\begin{proof}
Z definicji współrzędne $B_t^{(i)}/\sqrt t$ są niezależnymi
zmiennymi o rozkładzie $\Normal(0,1)$. Zatem
$|B_t|^2/t=\sum_{i=1}^d(B_t^{(i)}/\sqrt t)^2$,
co daje rozkład $\chi^2_d$ na mocy twierdzenia
o sumie kwadratów niezależnych zmiennych normalnych
z tomu I, rozdziału~\ref{I-ch:distributions}.
Każdy składnik kwadratowy ma średnią $1$ i wariancję $2$,
a niezależność pozwala dodać wartości oczekiwane
i wariancje.
\end{proof}

\begin{wniosek}[Gęstość promienia i jego momenty]
Dla $t>0$ zmienna $R_t=|B_t|$ ma na $(0,\infty)$
gęstość
\[
f_{R_t}(r)=
\frac{2^{1-d/2}}{\Gamma(d/2)t^{d/2}}
r^{d-1}e^{-r^2/(2t)}.
\]
Dla rzeczywistego $q$ moment $\E R_t^q$ jest
skończony wtedy i tylko wtedy, gdy $q>-d$;
w tym przypadku
\[
\E R_t^q=(2t)^{q/2}
\frac{\Gamma((d+q)/2)}{\Gamma(d/2)}.
\]
\end{wniosek}
\begin{proof}
Zmienna $Z=R_t^2/t$ ma rozkład $\chi_d^2$
o gęstości
$z^{d/2-1}e^{-z/2}/(2^{d/2}\Gamma(d/2))$ dla $z>0$
(por.~\eqref{I-eq:chisq}).
Podstawienie $z=r^2/t$ daje czynnik Jacobiego
$2r/t$ i wskazaną gęstość. Nie ma atomu w zerze,
ponieważ $B_t$ ma gęstość w $\R^d$.
W całce na moment podstawiamy $u=r^2/(2t)$:
\[
\E R_t^q=\frac{(2t)^{q/2}}{\Gamma(d/2)}
\int_0^\infty u^{(d+q)/2-1}e^{-u}\,\mathrm d\lambda_1(u).
\]
Dla każdego $q$ całka jest skończona przy
nieskończoności dzięki czynnikowi wykładniczemu.
Przy zerze jest skończona dokładnie wtedy,
gdy $(d+q)/2>0$. Otrzymujemy wzór z definicji
funkcji gamma.
\end{proof}

\begin{zadanie}[Oszacowanie prawdopodobieństwa wyjścia z sześcianu]
Niech $B$ będzie standardowym ruchem Browna w $\R^d$, $a,T>0$.
Oszacować $\Prob(\exists t\in[0,T]:\max_i|B_t^{(i)}|\ge a)$.

\emph{Rozwiązanie.} Jeśli proces opuszcza sześcian $(-a,a)^d$,
to dla pewnej współrzędnej $i$ zachodzi
$\sup_{t\le T} B_t^{(i)}\ge a$ lub
$\inf_{t\le T} B_t^{(i)}\le-a$.
Z nierówności sumy i zasady odbicia,
\[
\Prob\bigl(\sup_{t\le T}\|B_t\|_\infty\ge a\bigr)
\le4d\bigl(1-\Phi(a/\sqrt T)\bigr).
\]
Możemy zastąpić prawą stronę jej minimum z $1$.
Oszacowanie jest proste, choć nie musi być ostre,
ponieważ wielokrotnie zlicza trajektorie przekraczające
kilka ścian sześcianu. Z niezależności całych procesów
współrzędnych otrzymujemy również dokładną redukcję
do problemu jednowymiarowego:
\[
\Prob\left(\sup_{t\le T}\norm{B_t}_\infty\ge a\right)
=1-(1-p)^d,\qquad
p=\Prob\left(\sup_{t\le T}|W_t|\ge a\right).
\]
Zdarzenia pozostawania poszczególnych współrzędnych
w $(-a,a)$ przez cały $[0,T]$ są niezależne,
a każde z nich ma prawdopodobieństwo $1-p$.
\end{zadanie}

\section{Schematyczny rysunek procesu dwuwymiarowego}
\begin{center}
\begin{tikzpicture}[x=0.85cm,y=0.85cm,>=Latex]
 \draw[->,gray] (-3.2,0)--(3.7,0) node[right] {$x_1$};
 \draw[->,gray] (0,-2.5)--(0,2.7) node[above] {$x_2$};
 \draw[blue!70!black,very thick,rounded corners=1pt]
 (0,0)--(.15,.3)--(.45,.1)--(.1,.6)--(.3,.9)--(-.1,1.15)--(-.7,.9)
 --(-.55,.4)--(-.9,.2)--(-.4,-.2)--(-.8,-.7)--(-1.2,-.5)
 --(-1.45,-.85)--(-1.1,-1.1)--(-1.7,-1.4)--(-1.3,-1.7)
 --(-2.1,-1.5)--(-1.85,-1.1)--(-2.25,-.65)--(-1.9,-.25)
 --(-2.6,.2)--(-2.15,.6)--(-2.55,1.05)--(-2.1,1.3);
 \fill[red!75!black] (0,0) circle (2pt) node[below right] {$B_0$};
 \fill[blue!70!black] (-2.1,1.3) circle (2pt) node[above] {$B_T$};
 \draw[->,blue!70!black,thick] (-2.55,1.05)--(-2.1,1.3);
\end{tikzpicture}
\par\smallskip\small Schemat dwóch współrzędnych ruchu Browna w płaszczyźnie.
Łamana jest jedynie ilustracją: niemal każda rzeczywista trajektoria
jest ciągła i nieróżniczkowalna.
\end{center}


% ==============================================================
% ROZDZIAŁ 9. PROCESY ODNOWY
% ==============================================================
